% !TEX program = xelatex
% 論文① 完成に向けた実験計画（改訂版） 2026-09-29
\documentclass[11pt,a4paper]{article}

\usepackage[a4paper,margin=20mm]{geometry}
\usepackage{xeCJK}
\usepackage{fontspec}
\usepackage{amsmath,amssymb}
\usepackage{booktabs}
\usepackage{tabularx}
\usepackage{array}
\usepackage{xcolor}
\usepackage{enumitem}
\usepackage{fancyhdr}
\usepackage[hidelinks]{hyperref}

\setCJKmainfont{Hiragino Mincho ProN}
\setCJKsansfont{Hiragino Sans}
\setCJKmonofont{Hiragino Sans}
\setmonofont{Menlo}[Scale=0.80]
\xeCJKDeclareCharClass{CJK}{"2460 -> "24FF, "2500 -> "257F, "25A0 -> "25FF, "3000 -> "303F}

\newcolumntype{Y}{>{\raggedright\arraybackslash}X}
\setlist[itemize]{leftmargin=1.4em,itemsep=2pt,topsep=3pt}
\setlist[enumerate]{leftmargin=1.6em,itemsep=2pt,topsep=3pt}

\pagestyle{fancy}\fancyhf{}
\lhead{\small 論文① 完成に向けた実験計画（改訂版）}
\rhead{\small \thepage}
\renewcommand{\headrulewidth}{0.4pt}
\setlength{\headheight}{13pt}
\renewcommand{\tablename}{表}
\renewcommand{\abstractname}{要旨}

\title{\vspace{-12mm}\bfseries 論文①「自然言語から物理動作へ」\\[2mm]
\large 完成に向けた実験計画（改訂版）}
\author{法政大学 データサイエンスセンター長\quad 児玉 靖司}
\date{2026年9月29日}

\begin{document}
\maketitle
\thispagestyle{fancy}

\begin{abstract}
\noindent
2026年9月24日の計画書を、9月29日までの実機実験の結果に合わせて改訂する。
実機（DOFBOT + Xinu）は、同じAIPLプログラムで模型と実機の両方を動かし、
色の異なる立方体を挟んで置き場へ運ぶ作業を2回完遂した。
しかし成功率はまだ測っておらず、主実験（事前選別なしA／ありB の対応比較）に必要な
実機側のActor契約、投入ゲート、実行履歴の記録形式、採点器は未着手である。
本書は、残る5週間半（9月30日〜11月6日）で主実験を成立させるための実験の並びと日程、
実機の稼働時間の見積もり、撤退線、および先生に決めていただく事項を示す。
数値のうち実測したものと見積もりは区別して書く。
\end{abstract}

\section{到達点（2026年9月29日）}

\begin{table}[h]\centering\small
\caption{論文①の研究課題ごとの到達点}
\begin{tabularx}{\textwidth}{@{}lYl@{}}
\toprule
項目 & 状態 & 論文 \\
\midrule
RQ1 生成 & 未見の指示30件で予備実験済み（生成プログラム27件） & 7.1節 \\
RQ3 検証段 & 運動学（Tier-1）と剛体接触（Tier-2, MuJoCo）で27件を実行。Tier-1の誤った合格4件を検出 & 7.2節 \\
RQ2 両実行 & 同じ手順を模型と実機で実行。応答の欠落22.5\%、狭い可達域 & 7.3節 \\
実機の把持 & 色で振り分ける作業を20回実行し、最後の2回（青・緑）で完遂。模型が見逃した差を8つ同定 & 7.4節 \\
RQ4 主実験 & \textbf{未実施}（契約・ゲート・記録・採点器が無い） & --- \\
\bottomrule
\end{tabularx}
\end{table}

\paragraph{実機で確かめた道具。}次は実機で動いており、主実験の部品として使える。
\begin{itemize}
\item 手首カメラで立方体を探す手順（撮る→色の塊の重心を卓上へ逆射影→向け直す）。
\item 指の軸を合わせて段階的に下ろす手順（撮る姿勢と挟む姿勢を分ける）。
\item \textbf{挟めたかの判定}：指令135で閉じた後のサーボ6の読み戻しが133以下なら挟んでいる。
7回判定し、目視と矛盾した判定は無かった（挟めた2回はいずれも目視で確認）。ただし余裕は0（2回とも133）。
\item 寸法の一元管理（\texttt{geometry.json}、出典つき）。CG、逆運動学、逆射影が同じ値を読む。
\item 実行中の手を表示するPLAN（模型・実機・カメラ係の命令を一つの列に）。
\end{itemize}

\section{9月24日の計画からの遅れ}

\begin{table}[h]\centering\small
\caption{旧計画の週と実際}
\begin{tabularx}{\textwidth}{@{}lYY@{}}
\toprule
旧計画 & 予定 & 実際（9月29日） \\
\midrule
W0（9/24--27） & 既知座標で把持20回、成功率$\ge0.8$ & 把持は2回成功。\textbf{成功率は未測定} \\
W1（9/28--10/4） & 共通Actor契約（実機側）、投入ゲート、記録形式、E0 & \textbf{未着手}。代わりに寸法・指・カメラの実測と探索手順を作った \\
W2以降 & Tier-2、課題集合、採点器 & Tier-2は前倒しで完了。課題集合・採点器は未着手 \\
\bottomrule
\end{tabularx}
\end{table}

\noindent
W0が4日遅れ、W1が手つかずである。残りは5週間半であり、主実験の週（10/19--25）を動かさないために、
W1とW2の内容を圧縮して再配置する（第6節）。

\section{今回分かったことが計画に与える影響}

\begin{enumerate}
\item \textbf{可達域は土台の軸から14--23\,cmの帯に限られる}（指を8.3\,cmとした場合）。
課題の配置は、この帯の中に置き場と物体を並べる座標表として作る。
三段目の積み重ねは机上のどこでも届かない（7.3節）。
\emph{届かない指示}を課題集合に意図的に含めれば、Tier-1の正しい却下を測れる。
\item \textbf{Tier-2の模型が実機の実測に合っていない}。Tier-2（\texttt{harness/tier2}）は、把持点を掌から6\,cm、指の開きを約4.2\,cm（開き切り）から約0.2\,cmの直動とし、
立方体を1辺3.1\,cm、手首カメラを指の横（縦の画角56°）に置いて作られている。
実機では指は約8.3\,cm（推定）、開きは7段の実測表（開き切りで6.2\,cm、指令135で2.8\,cm）、立方体は3\,cm、
カメラは手首の上面にある。CGと色判定の係の模型も、Tier-2とはカメラの位置が異なる。
このままではシミュレーションの判定が実機の成否を予測するはずがない。
\textbf{Tier-2を\texttt{geometry.json}から組み立て直す}ことを主実験の前提とする。
\item \textbf{論文の記述に誤りが一つ残っている}。7.2節はTier-2のリンク長を「実機の実測値」としているが、
公開URDFの値である（7.4節と装置の段落は9月29日に訂正済み）。投稿前に直す。
\item \textbf{腕の基板が数分の動作で止まる}（4回）。Xinu側の自動リセットが引き金である疑いが濃い。
Mac側の読み取り制限の後の2回の実行では止まっていないが、2回では確かめたことにならない。
主実験の240回前後の実機実行に耐えるには、\texttt{system/arm.c}の自動リセットを外して焼き直す必要がある。
\item \textbf{挟めたかの判定には余裕が無い}。成功2回の読み戻しがしきい値と等しい。
しきい値を134に上げるか、挟む指令を立方体の幅に合わせて選び直し、成功と失敗の読み戻しの分布を30回で測る。
\item \textbf{卓上の他の色の物を立方体と取り違える}。主実験の机上には、立方体と置き場以外の色の物を置かない。
これは運用の規則として論文に書く。
\item \textbf{ブラウザが板のHTTPを止める}。主実験中はシミュレータの画面を開かないか、
板へ直接つながない構成にする（カメラ画像は係が取得する）。
\end{enumerate}

\section{実験の一覧（改訂）}

\begin{table}[h]\centering\small
\caption{実験の並び。E0$'$を新設し、主実験の前に置く}
\begin{tabularx}{\textwidth}{@{}clYl@{}}
\toprule
& 名称 & 内容 & 状態 \\
\midrule
E0$'$ & 実機の基礎性能 & 既知座標の立方体（青15・緑15、帯の中の5位置）を\texttt{sort.aipl}で30回。把持成功率、挟めた判定の読み戻しの分布、腕の基板の停止回数 & 新設 \\
E0 & 計器の検証 & 採点器$Y$と目視の一致（成功15・失敗15の場面）。投入ゲート$H$が範囲外の指令を止めること & 未着手 \\
E1 & 主実験A/B & 同一の初回生成プログラムを、事前選別なし(A)／あり(B)で対応づけて実機実行 & 未着手 \\
E2 & 生成の正しさ & 60指示で通過率（予備は30件で済み） & 一部済 \\
E3 & 却下の分類 & E1のデータから、段（静的・Tier-1・Tier-2）ごとの却下と誤った却下 & E1待ち \\
E4 & 誤った合格 & E1のデータから、Bが合格させ実機で失敗した事例と原因 & E1待ち \\
E6 & 人手確認済み & 正しいと確認したプログラム（\texttt{sort.aipl}など）を両バックエンドで実行 & 一部済 \\
E5 & 修正ループ & 余力があれば & 保留 \\
\bottomrule
\end{tabularx}
\end{table}

\paragraph{E0$'$の合格条件。}把持成功率$\ge0.8$（30回中24回以上）、かつ30回の間に腕の基板が止まらないこと。
成功率が0.5未満なら第7節の撤退線に移る。0.5以上0.8未満なら、失敗を分類して一つだけ直し、もう30回行う。
\textbf{結果を見てから合格条件を変えない。}

\section{試行数と実機の稼働時間}

検出力の計算（9月24日の計画書）は変えない。McNemar正確検定で検出力0.8を得るには、
効果0.25で62件が要る。\textbf{異なる指示60件×配置2通り}を計画値とし、指示単位のクラスタ・ブートストラップで解析する。

\begin{table}[h]\centering\small
\caption{実機の実行回数と稼働時間の見積もり（1回3分は見積もりであり実測ではない）}
\begin{tabularx}{\textwidth}{@{}lrrY@{}}
\toprule
実験 & 実機の実行 & 時間 & 備考 \\
\midrule
E0$'$ & 30--60 & 1.5--3\,h & 追加の30回を含む \\
E0 & 30 & 1.5\,h & 採点器の一致確認 \\
予備（10指示×2配置×A/B） & $\le$40 & $\le$2\,h & 本実験に入れない \\
E1（60指示×2配置） & A 120 + B $\le$120 & $\le$12\,h & 配置の戻しと較正を含まない \\
\midrule
計 & $\le$250 & $\le$18.5\,h & 3時間の実機セッション7回前後 \\
\bottomrule
\end{tabularx}
\end{table}

\noindent
各セッションには、安全確認と目視判定のために操作者が一人つく。
9月28--29日の実行では、1回の実行で試行を6回繰り返した例があり（20回目）、そのときは3分を超えている。
E0$'$で1回の実行時間を実測し、この表を引き直す。

\section{日程（9月30日〜11月6日）}

\subsection*{W1：9月30日（火）〜10月4日（日）— 実機を主実験に耐える状態にする}
\begin{enumerate}
\item Xinu：\texttt{system/arm.c}の自動リセットを外す。あわせて\texttt{d4eac11}で消えたファンの温度追従と、
動かない起動時の仕事（アクター載せ・カメラ開始）を直す。焼いて、冷起動3回で\texttt{/version}とHTTPを確認する。
\item 仮置きの寸法を物差しで測る：指の長さ、手首からカメラ中心まで、指の軸からカメラ中心まで、台座の高さ。
画角は、既知の大きさの立方体を既知の距離で撮って求める。\texttt{geometry.json}を「実測」に更新する。
\item Tier-2（MuJoCo）を\texttt{geometry.json}から組み立て直し、指の開きを実測表に合わせる。
\item \textbf{E0$'$を実施する}（30回）。挟めた判定のしきい値を読み戻しの分布から決める。
\end{enumerate}

\subsection*{W2：10月5日（月）〜10月11日（日）— 契約・ゲート・記録・採点器}
\begin{enumerate}
\item 実機側のActor契約を、今ある部品の上に実装する：\texttt{observe}は探索、\texttt{pick}は下ろしと挟めた判定、
\texttt{place}は置き場への下ろし。結果は型付き（\texttt{OK, NOT\_FOUND, UNREACHABLE, NOT\_HELD, TIMEOUT}）。
要求IDと期限を付け、\textbf{時間切れを成功と扱わない}。
\item 投入ゲート$H$：関節範囲、卓上の最低高さ、可達帯（14--23\,cm）。両条件で有効。
\item 実行履歴をJSONLで残す（時刻、プログラムID、呼出し、引数、結果、場面の版）。PLANの列をそのまま使う。
\item 採点器$Y$：置き場を撮って色と位置を判定する。\textbf{E0を実施する}。
\end{enumerate}

\subsection*{W3：10月12日（月）〜10月18日（日）— 課題集合と予備実験}
\begin{enumerate}
\item 課題集合60件を六群で確定し、例示用と未見用を\emph{プロンプトを書く前に}分ける。
可達帯の外や三段目を求める指示を群ごとに含める。
\item 予備実験（10指示）を端から端まで通す。予備のデータは本実験に入れない。
\item 検出力を予備データで引き直し、試行数を確定する。\textbf{10月18日に計画を凍結する}。
\end{enumerate}

\subsection*{W4：10月19日（月）〜10月25日（日）— 本実験}
課題順と条件順を無作為化する。物体と置き場は座標表どおりに毎回戻す。操作者の介入はすべて記録する。
条件Aで安全ゲートが止めた試行はAの失敗として数える。

\subsection*{W5：10月26日（月）〜11月1日（日）— 副次実験と解析}
E6を先に、E5は余力があれば。成功率差の信頼区間、McNemar、群別の結果、所要時間の分布、失敗の分類。

\subsection*{W6：11月2日（月）〜11月6日（金）— 執筆}
結果の節を書き、要旨・考察・結論を結果に合わせて改稿する。7.2節のリンク長の記述を直す。
機体・制御器・カメラ・グリッパー・シミュレータ・LLMとAIPL処理系の版・寸法表を明記する。

\section{リスクと撤退線}

\paragraph{E0$'$で成功率が0.5未満。}(1)挟む指令と高さを読み戻しの分布から選び直す。
(2)立方体を指の開きに合う大きさ（3.5--4\,cm）に替える。(3)課題を「押して移動」に変える。
いずれも、採用した対処を論文に明記する。

\paragraph{腕の基板が焼き直し後も止まる。}自動リセットが原因ではなかったことになる。
電源（アダプタの容量）を疑い、別のアダプタで30回を試す。主実験のセッションを短く区切り、
止まった試行は「制御器の失敗」として分類して分母に含める。

\paragraph{W2の実装が遅れる。}採点器$Y$を目視判定で代替し、その旨を明記する（E0は不要になる）。
主実験の週は動かさない。

\paragraph{効果が小さいとW3で分かる。}群を減らして指示数を増やすか、差が小さいことを区間で示す主張に変える。
結果を見てから試行数を足さない。

\section{先生に決めていただく事項}

\begin{enumerate}
\item \textbf{物体の位置の入力}：主実験で、手首カメラによる探索を使うか、座標表の値を入力するか。
後者は知覚の誤りを比較から除けるが、論文の射程が狭まる。どちらも両条件で同じ入力にする。
\item \textbf{Code as Policies型との比較}（論文の副次比較）を残すか。W2--W3の工数では難しい。削るなら主張から外す。
\item \textbf{Xinuの焼き直し}（自動リセットの削除）をW1の初めに行ってよいか。
\item 机上の物：立方体の色と数、置き場の数（現状は青・緑、左右の二つ）。
\item 主実験のセッションに立ち会える日（3時間×7回前後）。
\end{enumerate}

\section*{決定事項（2026年9月29日、先生の回答）}
\begin{enumerate}
\item 物体の位置は\textbf{手首カメラで探す}（座標表の入力は使わない）。両条件で同じ探索を用いる。
\item Code as Policies型との比較は\textbf{論文に残す}。W2--W3に実装の工数を確保する。
\item Xinuは\textbf{焼き直す}。9月29日に自動リセットの削除とファンの温度追従の復旧を焼いた（\texttt{Sep 29 2026 09:57:23}）。
\item 次の実験では\textbf{緑と青の2つを同時に置き、1回の実行で両方を運ぶ}。
\item 立会いは\textbf{いつでも可}。
\end{enumerate}

\end{document}
